\begin{ouverture}[geometrie]{13}{Droites, distances et cercles}{Partie 3 · Espace et géométrie}
{Notations · Distance · Milieu · Cercle et disque}
\objectif{Employer un vocabulaire précis et coder les propriétés utiles.}
\objectif{Utiliser la distance entre deux points et le milieu d’un segment.}
\objectif{Caractériser cercle et disque par une distance au centre.}
\prerequis{Règle, équerre et compas ; perpendicularité et parallélisme.}
\end{ouverture}
\pagestyle{otMain}\essentiel\fontsize{11}{15}\selectfont
\phantomsection\label{cours-13-1}
\essSection{1}{Objets, notations et positions de droites}
\begin{definition}Par deux points A et B passe une seule droite, notée (AB). Le segment [AB] est la partie de cette droite comprise entre A et B ; la demi-droite [AB) commence en A et passe par B. La longueur du segment [AB] se note AB, sans crochets.\end{definition}
\label{fig-cours-13-1-cours-objets}\figureGeo{cours-objets}

\begin{propriete}Deux droites perpendiculaires se coupent en formant un angle droit ; deux droites parallèles ne se coupent pas. On code un angle droit par un petit carré et des longueurs égales par un même signe.\end{propriete}
\essSection{2}{Distance et milieu d’un segment}
\begin{definition}La distance entre deux points A et B est la longueur AB du segment [AB]. Le milieu du segment [AB] est le point M de ce segment tel que $AM=MB$ ; alors $AM=MB=AB:2$.\end{definition}
\begin{savoirfaire}[Placer un milieu]Pour placer le milieu d’un segment de 7 cm, on place sur le segment le point situé à $7:2=3{,}5$ cm de l’une des extrémités, puis on code les deux moitiés.\end{savoirfaire}
\clearpage
\phantomsection\label{cours-13-2}
\essSection{3}{Cercle, disque et éléments associés}
\begin{definition}Le cercle de centre O et de rayon $r$ est formé de tous les points situés à la distance $r$ du point O. Le disque de centre O et de rayon $r$ est formé de tous les points dont la distance à O est inférieure ou égale à $r$ : c’est le cercle et son intérieur.\end{definition}
\label{fig-cours-13-2-cours-cercle}\figureGeo{cours-cercle}

\begin{definition}Un rayon joint le centre à un point du cercle. Une corde joint deux points du cercle ; une corde qui passe par le centre est un diamètre. Un diamètre mesure deux rayons.\end{definition}
\essSection{4}{Décrire des conditions de distance}
\begin{savoirfaire}[Comparer à un rayon]Pour situer un point M par rapport au cercle de centre O et de rayon 3 cm, on compare OM à 3 cm : si $OM=3$ cm, M est sur le cercle ; si $OM<3$ cm, M est à l’intérieur du cercle ; si $OM>3$ cm, M est à l’extérieur.\end{savoirfaire}
\begin{exemple}Les points situés à 2 cm de A et à 3 cm de B sont à la fois sur le cercle de centre A de rayon 2 cm et sur le cercle de centre B de rayon 3 cm : ce sont leurs points d’intersection.\end{exemple}

\clearpage\rubrique{Découverte}
\begin{blocLivre}{ex-13-1}\begin{exercice}[1]{Nommer les objets}
Nommer chacun des trois objets tracés avec la bonne notation.
\label{fig-ex-13-1}\figureGeo{dis-1}
\end{exercice}\end{blocLivre}
\begin{blocLivre}{ex-13-2}\begin{exercice}[1]{Lire les codages}
Écrire toutes les informations données par les codages du quadrilatère ABCD.
\label{fig-ex-13-2}\figureGeo{dis-2}
\end{exercice}\end{blocLivre}
\begin{blocLivre}{ex-13-3}\begin{exercice}[1]{Le milieu}
Le segment [AB] mesure 7 cm et M est son milieu. Calculer AM. Le point N du segment [AB], placé à 3 cm de A, est-il le milieu de [AB] ?
\label{fig-ex-13-3}\figureGeo{dis-3}
\end{exercice}\end{blocLivre}
\begin{blocLivre}{ex-13-4}\begin{exercice}[1]{Vocabulaire du cercle}
Pour chaque segment tracé, dire s’il s’agit d’un rayon, d’un diamètre ou d’une corde.
\label{fig-ex-13-4}\figureGeo{dis-4}
\end{exercice}\end{blocLivre}
\begin{blocLivre}{ex-13-5}\begin{exercice}[1]{Sur le cercle ou non ?}
Le cercle a pour centre O et pour rayon 3 cm. On sait que $OM=3$ cm, $ON=2{,}4$ cm et $OP=3{,}6$ cm. Situer M, N et P par rapport au cercle.
\label{fig-ex-13-5}\figureGeo{dis-5}
\end{exercice}\end{blocLivre}

\rubrique{Entraînement}
\begin{blocLivre}{ex-13-6}\begin{exercice}[2]{Rayon et diamètre}
Le cercle $\mathcal C_1$ a pour diamètre $4{,}6$ cm et le cercle $\mathcal C_2$ a pour rayon $1{,}4$ cm. Calculer le rayon de $\mathcal C_1$ et le diamètre de $\mathcal C_2$. Lequel est le plus grand cercle ?
\label{fig-ex-13-6}\figureGeo{dis-6}
\end{exercice}\end{blocLivre}
\begin{blocLivre}{ex-13-7}\begin{exercice}[2]{À égale distance}
Les points A et B sont tels que $AB=4$ cm. Construire tous les points situés à 2 cm de A et à 3 cm de B. Combien y en a-t-il ?
\label{fig-ex-13-7}\figureGeo{dis-7}
\end{exercice}\end{blocLivre}
\begin{blocLivre}{ex-13-8}\begin{exercice}[2]{Un cercle de diamètre donné}
Le segment [AB] mesure 6 cm. Placer son milieu M, puis tracer le cercle de diamètre [AB]. Quel est son rayon ?
\label{fig-ex-13-8}\figureGeo{dis-8}
\end{exercice}\end{blocLivre}
\begin{blocLivre}{ex-13-9}\begin{exercice}[2]{La corde la plus longue}
Dans le cercle de centre O et de rayon $2{,}5$ cm, trois cordes sont tracées. Laquelle est la plus longue ? Peut-on tracer une corde de 6 cm dans ce cercle ?
\label{fig-ex-13-9}\figureGeo{dis-9}
\end{exercice}\end{blocLivre}
\begin{blocLivre}{ex-13-10}\begin{exercice}[2]{Comparer au compas}
Sans règle graduée, ranger les segments [AB], [CD] et [EF] du plus court au plus long, à l’aide du compas.
\label{fig-ex-13-10}\figureGeo{dis-10}
\end{exercice}\end{blocLivre}

\rubrique{Problèmes}
\begin{blocLivre}{ex-13-11}\begin{exercice}[3]{L’antenne}
Sur la carte, 1 cm représente 1 km. Une antenne A capte les téléphones jusqu’à 3 km. Quels villages sont couverts ? Les distances à l’antenne sont : V1 $2{,}1$ km ; V2 $3{,}4$ km ; V3 $2{,}9$ km ; V4 3 km.
\label{fig-ex-13-11}\figureGeo{dis-11}
\end{exercice}\end{blocLivre}
\begin{blocLivre}{ex-13-12}\begin{exercice}[3]{Le trésor}
Sur ce plan, 1 cm représente 1 m. Un trésor est enterré à 4 m du pin P et à 3 m du rocher R, qui sont distants de 5 m. Où peut-il se trouver ?
\label{fig-ex-13-12}\figureGeo{dis-12}
\end{exercice}\end{blocLivre}
\begin{blocLivre}{ex-13-13}\begin{exercice}[3]{La chèvre}
Une chèvre est attachée au piquet O par une corde de 5 m. Sur le plan, 1 cm représente 2 m. Les fleurs sont à 4 m (F1), $5{,}5$ m (F2) et 5 m (F3) du piquet. Lesquelles peut-elle brouter ?
\label{fig-ex-13-13}\figureGeo{dis-13}
\end{exercice}\end{blocLivre}
\begin{blocLivre}{ex-13-14}\begin{exercice}[3]{Deux cercles qui se touchent}
Les cercles de centres A et B ont le même rayon, 3 cm, et $AB=6$ cm. Ils se touchent en un seul point M, situé sur le segment [AB]. Expliquer pourquoi M est le milieu de [AB].
\label{fig-ex-13-14}\figureGeo{dis-14}
\end{exercice}\end{blocLivre}
\begin{blocLivre}{ex-13-15}\begin{exercice}[3]{Vrai ou faux}
Yanis affirme : « Si $MA=MB$, alors M est le milieu du segment [AB]. » Le point M de la figure vérifie $MA=MB=3$ cm et $AB=4$ cm. Yanis a-t-il raison ?
\label{fig-ex-13-15}\figureGeo{dis-15}
\end{exercice}\end{blocLivre}

\clearpage
